\chapter{Results}
\label{chap:results}

\textbf{This chapter is empty. No result exists yet.}

At the time of this revision the repository contains no capture from the TV2
session. The most recent non-simulated stereo pairs in
\texttt{data/captures/} date from February 2026, and the most recent
simulated ones from September 2026. The record sheet of the pre-session
parameter artifact has not been filled in.

Nothing in this chapter should be written before the session. What follows is
the structure to fill, so that the writing afterwards is transcription rather
than reconstruction.

\section{Session record}

Transcribe Table \ref{tab:record} here, completed. Every \emph{not measured}
entry keeps its reason.

\section{Calibration outcome}

Report against Table \ref{tab:crit-calib}: the accepted view count per tilt
bin, the relative standard deviations of $f_x$ and $f_y$, the standard
deviation of each distortion coefficient against its own magnitude, the
principal point offset, and whether the tool reported the session complete.

Report the reprojection RMS as well, but report it \emph{after} the
uncertainties and label it as the training residual it is. If A1.2 was not
met, state the value achieved and carry the implied scale error forward into
every length quoted in this chapter.

\section{Illumination outcome}

Report against Table \ref{tab:crit-light}, one row per criterion, with the
measured value and pass or fail. Where a criterion failed, state what was
tried.

\section{Pipeline outcome}

Report against Table \ref{tab:crit-pipeline}, aggregated across the analysed
pairs, with the per-set breakdown from Table \ref{tab:capture-plan}. Every
arc length is reported with its retained sample fraction beside it.

\section{The two anticipated failure modes}

\textbf{Parametric matching, from set S8.} The measured discrepancy against
the tape-measured length, with its sign, and how it compares to the S1 pairs
of the same tether.

\textbf{Rolling-shutter shear, from sets S4 and S6.} A qualitative
description of whether shear was visible at all, and at what point in the
motion the tracker or the matching visibly stopped working. No angular rate
is quoted, because none was measured.

\section{Verdict}

State the overall verdict using the definitions at the end of Chapter
\ref{chap:criteria}: pass, qualified pass, or fail. State it before
discussing what it means.

\section{What the next session should change}

Written last, from what actually went wrong rather than from what was
anticipated here.

\todo[inline]{Pending: run the session and complete this chapter. Until it is
complete, this deliverable is a protocol, not a report, and should not be
cited as evidence that the instrument has been validated.}
